% Lecture 11 notes: \input by notes.tex, \subimport'ed by the semester book.
% Exam questions (экзамен): 25, 26, 28 (lattices, Plan 4); the lecture follows
% HPS ch. 7: congruential cryptosystem, lattices and bases, det and Hadamard ratio,
% SVP/CVP, Hermite and Minkowski, Gauss reduction, Babai's algorithms.
\begin{lecture}{11}{Решётки. Алгоритмы Гаусса и Бабаи}

\section{Зачем криптографу решётки}

Все криптосистемы, которые мы изучали до сих пор (Диффи"--~Хеллман, RSA, Эль-Гамаль, протоколы
на эллиптических кривых), опираются на две задачи: дискретное логарифмирование и факторизацию.
Обе решаются за полиномиальное время на квантовом компьютере (алгоритм Шора), поэтому сегодня
активно используются криптосистемы, стойкость которых основана на \emph{задачах о решётках}:
найти кратчайший вектор решётки или ближайший к данной точке вектор решётки. В 2024~г.
американский институт стандартов NIST утвердил стандарты постквантовой криптографии ML-KEM
(Kyber) и ML-DSA (Dilithium) "--- обе схемы построены на решётках.

Решётки полезны и криптоаналитику: многие атаки (на рюкзачные криптосистемы, на RSA с малой
экспонентой, на подписи с «плохими» случайными числами) сводятся к поиску короткого вектора
в подходящей решётке. Начнём с самого простого примера, в котором видны обе стороны:
криптосистема, атака на которую "--- поиск кратчайшего вектора в решётке размерности~$2$.

\subsection{Конгруэнтная криптосистема}

Следующая игрушечная криптосистема \cite[разд.~7.1]{hps} "--- двумерный прообраз системы NTRU.

\begin{itemize}
  \item \textbf{Ключи (Алиса).} Выбирает большое число $q$ и целые $f, g$ с условиями
    \[
      f < \sqrt{q/2}, \qquad \sqrt{q/4} < g < \sqrt{q/2}, \qquad \gcd(f, qg) = 1,
    \]
    вычисляет $\congr{h}{f^{-1} g}{q}$, $0 < h < q$. Открытый ключ "--- $(q, h)$, закрытый "---
    $(f, g)$.
  \item \textbf{Шифрование (Боб).} Открытый текст "--- целое $m$, $0 < m < \sqrt{q/4}$. Боб
    выбирает случайное $r$, $0 < r < \sqrt{q/2}$, и отправляет $\congr{e}{rh + m}{q}$, $0 \le e < q$.
  \item \textbf{Расшифрование (Алиса).} Вычисляет $\congr{a}{fe}{q}$, $0 \le a < q$, затем
    $\congr{b}{f^{-1} a}{g}$, $0 \le b < g$. Тогда $b = m$.
\end{itemize}

\begin{figure}[ht]
\centering
\begin{tikzpicture}[x=1cm,y=1cm]
  \node (A) at (0,0) {\personicon{caBlue}{Алиса}};
  \node (B) at (9,0) {\personicon{caRed}{Боб}};
  \draw[very thick,caGray,-{Stealth[length=3mm]}] (1,0.55) -- (8,0.55)
    node[midway,above,font=\small] {$(q, h)$};
  \draw[very thick,caGray,-{Stealth[length=3mm]}] (8,0.0) -- (1,0.0)
    node[midway,below,font=\small] {$e \equiv rh + m$};
  \node[font=\small,align=left,anchor=north] at (0,-1.0)
    {$f, g$ малы, $h \equiv f^{-1}g$\\ $a \equiv fe \pmod q$\\ $m \equiv f^{-1}a \pmod g$};
  \node[font=\small,align=left,anchor=north] at (9,-1.0) {$m$, случайное $r$};
\end{tikzpicture}
\caption{Конгруэнтная криптосистема}
\label{l11:fig-congr}
\end{figure}

\begin{proposition}\label{l11:congr-correct}
Расшифрование конгруэнтной криптосистемы корректно: $b = m$.
\end{proposition}

\begin{proof}
По модулю $q$ имеем $a \equiv fe \equiv f(rh + m) \equiv rg + fm$, так как $fh \equiv g$. Из
ограничений на параметры
\[
  0 < rg + fm < \sqrt{q/2}\sqrt{q/2} + \sqrt{q/2}\sqrt{q/4} = \Bigl(\tfrac12 + \tfrac{1}{2\sqrt2}\Bigr) q < q,
\]
поэтому $a$ (вычет из $[0, q)$) \emph{равен} целому числу $rg + fm$, а не только сравним с ним.
Значит, $\congr{a}{fm}{g}$, и $\congr{b}{f^{-1}a}{g}$ равно $m \bmod g$. Наконец,
$0 < m < \sqrt{q/4} < g$, т.\,е. $b = m$.
\end{proof}

\begin{example}\label{l11:congr-example}
$q = 1\,000\,003$ (простое), $f = 325$, $g = 534$: $\sqrt{q/2} \approx 707{,}1$,
$\sqrt{q/4} \approx 500{,}0$, $\gcd(325, 534) = 1$. Открытый ключ $h \equiv f^{-1} g = 529\,234$.
Боб шифрует $m = 478$ со случайным $r = 353$: $e \equiv 353 \cdot 529\,234 + 478 \equiv 819\,522$.
Алиса: $a \equiv 325 \cdot 819\,522 \equiv 343\,852 \pmod q$; действительно,
$rg + fm = 188\,502 + 155\,350 = 343\,852$. Далее $f^{-1} \equiv 511 \pmod{534}$ и
$\congr{511 \cdot 343\,852}{478}{534}$.
\end{example}

\paragraph{Атака.} Ева знает $(q, h)$. Закрытый ключ удовлетворяет $fh \equiv g \pmod q$, т.\,е.
$fh = g + Rq$ при некотором целом $R$, и
\[
  f\,(1, h) - R\,(0, q) = (f, g).
\]
Следовательно, $(f, g)$ "--- целочисленная комбинация векторов $v_1 = (1, h)$, $v_2 = (0, q)$:
это вектор \emph{решётки}, порождённой $v_1, v_2$. Длины $v_1, v_2$ порядка $q$, а
$\norm{(f, g)} < \sqrt{q}$ "--- вектор $(f, g)$ очень короткий по сравнению с ними. Ниже мы увидим, что в
решётке размерности~$2$ кратчайший вектор находится быстро (алгоритм Гаусса,
разд.~\ref{l11:sec-gauss}), и почти всегда это $\pm(f, g)$ (или другой вектор, пригодный
для расшифрования). Значит, конгруэнтная криптосистема нестойка при \emph{любом} $q$: стойкость
нужно искать в решётках большой размерности.

\checkyourself{Почему в конгруэнтной криптосистеме нельзя взять $g$ намного меньше
$\sqrt{q/4}$? Где в доказательстве утверждения~\ref{l11:congr-correct} используется
условие $\gcd(f, qg) = 1$?}

\section{Решётки и их базисы}

\subsection{Векторные пространства: напоминание}

В $\RR^n$ используем скалярное произведение $v \cdot w = \sum v_i w_i$ и евклидову норму
$\norm{v} = \sqrt{v \cdot v}$; векторы записываем \emph{строками}. Для набора векторов
$v_1, \dots, v_n$ будем писать $B$ для матрицы со строками $v_1, \dots, v_n$.

\begin{definition}
\emph{Ортогонализация Грама"--~Шмидта} линейно независимых $v_1, \dots, v_n$:
\[
  v_1^* = v_1, \qquad v_i^* = v_i - \sum_{j < i} \mu_{ij} v_j^*, \qquad
  \mu_{ij} = \frac{v_i \cdot v_j^*}{\norm{v_j^*}^2} \quad (j < i).
\]
\end{definition}

Векторы $v_1^*, \dots, v_n^*$ попарно ортогональны, $\Span(v_1^*, \dots, v_i^*) = \Span(v_1, \dots,
v_i)$ при каждом $i$, и $v_i^*$ "--- проекция $v_i$ на ортогональное дополнение к
$\Span(v_1, \dots, v_{i-1})$. В матричном виде $B = M B^*$, где $M$ "--- нижнетреугольная матрица с
единицами на диагонали и элементами $\mu_{ij}$ под ней. Заметим, что $v_i^*$, вообще говоря, не
являются целочисленными комбинациями $v_j$.

\subsection{Решётка, базис, смена базиса}

\begin{definition}
Пусть $v_1, \dots, v_n \in \RR^m$ линейно независимы. \emph{Решёткой}, порождённой ими, называется
множество целочисленных линейных комбинаций
\[
  \Lat = \Lat(v_1, \dots, v_n) = \{a_1 v_1 + \cdots + a_n v_n : a_1, \dots, a_n \in \ZZ\}
  = \{xB : x \in \ZZ^n\}.
\]
Набор $v_1, \dots, v_n$ "--- \emph{базис} решётки, $n$ "--- её \emph{размерность}. Если
$\Lat \subseteq \ZZ^m$, решётка называется \emph{целочисленной}.
\end{definition}

Чаще всего $m = n$ (решётка полного ранга); в этой лекции, если не оговорено иное, $m = n$ и
$B$ "--- невырожденная квадратная матрица. Примеры: $\ZZ^n$ (базис "--- единичные векторы);
решётка $\{(x, y) \in \ZZ^2 : \congr{xh}{y}{q}\}$ из атаки на конгруэнтную криптосистему
(базис $(1, h), (0, q)$); решётка $\Lat((2, 1), (-1, 1))$ на рис.~\ref{l11:fig-bases}.
Решётка "--- аддитивная подгруппа $\RR^n$, причём \emph{дискретная}: в каждом шаре лишь
конечное число её точек (это следует из невырожденности $B$: координаты $x = vB^{-1}$ точек
шара ограничены). Верно и обратное "--- всякая дискретная подгруппа $\RR^n$ есть решётка
(без доказательства, см.~\cite[гл.~16]{galbraith}).

У решётки размерности $n \ge 2$ бесконечно много базисов. Какие именно?

\begin{theorem}[смена базиса]\label{l11:basis-change}
Пусть $B$ и $B'$ "--- матрицы базисов (по строкам). Решётки $\Lat(B)$ и $\Lat(B')$ совпадают тогда
и только тогда, когда $B' = UB$ для некоторой $U \in \mathrm{GL}_n(\ZZ)$, т.\,е. целочисленной
матрицы $U$ с $\det U = \pm 1$.
\end{theorem}

\begin{proof}
($\Leftarrow$) Строки $B' = UB$ "--- целочисленные комбинации строк $B$, поэтому
$\Lat(B') \subseteq \Lat(B)$. Обратная матрица $U^{-1} = \frac{1}{\det U}\operatorname{adj} U =
\pm \operatorname{adj} U$ тоже целочисленная (присоединённая матрица составлена из миноров
$U$), и $B = U^{-1} B'$ даёт $\Lat(B) \subseteq \Lat(B')$.

($\Rightarrow$) Каждая строка $B'$ лежит в $\Lat(B)$, т.\,е. $B' = UB$ с целочисленной $U$;
аналогично $B = VB'$ с целочисленной $V$. Тогда $B = VUB$, т.\,е. $(VU - I)B = 0$, и, так как
строки $B$ линейно независимы, $VU = I$. Значит, $\det V \cdot \det U = 1$, где оба
определителя "--- целые числа, откуда $\det U = \pm 1$.
\end{proof}

\begin{corollary}\label{l11:det-invariant}
Величина $\abs{\det B}$ не зависит от выбора базиса решётки.
\end{corollary}

\begin{proof}
$\det B' = \det U \det B = \pm \det B$.
\end{proof}

\begin{definition}
\emph{Определитель} (\emph{кообъём}) решётки $\Lat$ полного ранга: $\det \Lat = \abs{\det B}$
для любого базиса $B$. (Для решётки неполного ранга $\det \Lat = \sqrt{\det(BB^T)}$.)
\end{definition}

\begin{remark}
Равенство определителей \emph{не} означает совпадения решёток. Например, $\Lat((2, 0), (0, 1))$
и $\Lat((1, 0), (0, 2))$ имеют определитель $2$, но $(1, 0)$ лежит только во второй из них.
Чтобы проверить, порождают ли $B$ и $B'$ одну решётку, нужно найти $U = B' B^{-1}$ и проверить,
что она \emph{целочисленная} и $\det U = \pm 1$. Если $U$ целочисленная, но $\abs{\det U} = k > 1$,
то $\Lat(B')$ "--- подрешётка $\Lat(B)$ индекса $k$.
\end{remark}

\taskexample{lattice-det-hadamard}

\subsection{Фундаментальная область}

\begin{definition}
\emph{Фундаментальная область} (фундаментальный параллелепипед) базиса $v_1, \dots, v_n$:
\[
  \mathcal{F}(v_1, \dots, v_n) = \{t_1 v_1 + \cdots + t_n v_n : 0 \le t_i < 1\}.
\]
\end{definition}

\begin{proposition}\label{l11:tiling}
Каждый вектор $w \in \RR^n$ единственным образом представляется в виде $w = t + v$, где
$t \in \mathcal F$, $v \in \Lat$. Иначе говоря, сдвиги $\mathcal F + v$, $v \in \Lat$, покрывают
$\RR^n$ без пересечений.
\end{proposition}

\begin{proof}
Так как $v_1, \dots, v_n$ "--- базис $\RR^n$, $w = \sum \alpha_i v_i$ с вещественными $\alpha_i$.
Положим $v = \sum \floor{\alpha_i} v_i \in \Lat$ и $t = \sum \{\alpha_i\} v_i \in \mathcal F$
($\{\cdot\}$ "--- дробная часть). Если $t + v = t' + v'$, то коэффициенты $t$ и $t'$ в базисе $v_i$
отличаются на целые числа и лежат в $[0, 1)$, поэтому совпадают; значит, $t = t'$ и $v = v'$.
\end{proof}

\begin{theorem}\label{l11:volume}
$\operatorname{vol}(\mathcal F) = \abs{\det B} = \det \Lat$. В частности, объём
фундаментальной области одинаков для всех базисов решётки.
\end{theorem}

\begin{proof}
$\mathcal F$ "--- образ единичного куба $[0, 1)^n$ при линейном отображении $x \mapsto xB$, а
линейное отображение умножает объёмы на модуль своего определителя (формула замены
переменных в кратном интеграле). Объём куба равен $1$.
\end{proof}

Таким образом, $\det \Lat$ "--- это «объём, приходящийся на одну точку решётки»: в большом шаре
радиуса $R$ лежит примерно $\operatorname{vol}(\text{шара})/\det \Lat$ точек решётки.

\begin{figure}[ht]
\centering
\latticefigure[scale=0.85]{2}{1}{-1}{1}\qquad
\latticefigure[scale=0.85]{-4}{1}{-5}{2}
\caption{Одна и та же решётка $\det \Lat = 3$ в хорошем базисе $v_1 = (2, 1)$, $v_2 = (-1, 1)$,
$\mathcal H \approx 0{,}974$ (слева) и в плохом базисе $v_1 = (-4, 1)$, $v_2 = (-5, 2)$,
$\mathcal H \approx 0{,}368$ (справа). Площади закрашенных параллелограммов равны~$3$.}
\label{l11:fig-bases}
\end{figure}

\subsection{Неравенство Адамара}

\begin{theorem}[неравенство Адамара]\label{l11:hadamard}
Для любого базиса $v_1, \dots, v_n$ решётки $\Lat$
\[
  \det \Lat = \prod_{i=1}^n \norm{v_i^*} \le \prod_{i=1}^n \norm{v_i},
\]
причём равенство достигается тогда и только тогда, когда векторы $v_i$ попарно ортогональны.
\end{theorem}

\begin{proof}
Так как $B = MB^*$ и $\det M = 1$, имеем $\det B = \det B^*$. Строки $B^*$ попарно
ортогональны, поэтому $B^*(B^*)^T = \operatorname{diag}(\norm{v_1^*}^2, \dots, \norm{v_n^*}^2)$ и
$(\det B^*)^2 = \prod \norm{v_i^*}^2$. По теореме Пифагора
$\norm{v_i}^2 = \norm{v_i^*}^2 + \sum_{j<i} \mu_{ij}^2 \norm{v_j^*}^2 \ge \norm{v_i^*}^2$,
с равенством при всех $i$ тогда и только тогда, когда все $\mu_{ij} = 0$, т.\,е. когда $v_i$
попарно ортогональны.
\end{proof}

\begin{definition}
\emph{Коэффициент Адамара} базиса $B = (v_1, \dots, v_n)$:
\[
  \mathcal H(B) = \left(\frac{\det \Lat}{\norm{v_1}\norm{v_2}\cdots\norm{v_n}}\right)^{1/n}
  \in (0, 1]
\]
(в задачах он обозначается также $\HR$). Чем ближе $\mathcal H(B)$ к $1$, тем «ортогональнее»
базис; базис с $\mathcal H(B)$, близким к $1$, называют \emph{хорошим}, а с малым $\mathcal H(B)$ "---
\emph{плохим}.
\end{definition}

В размерности $2$: $\mathcal H(B)^2 = \abs{\sin\varphi}$, где $\varphi$ "--- угол между $v_1$ и
$v_2$. Для рис.~\ref{l11:fig-bases}:
$\mathcal H = \bigl(3/(\sqrt5 \sqrt2)\bigr)^{1/2} \approx 0{,}974$ и
$\mathcal H = \bigl(3/(\sqrt{17}\sqrt{29})\bigr)^{1/2} \approx 0{,}368$.

\checkyourself{Порождают ли векторы $(3, 1), (1, 2)$ и $(4, 3), (2, 4)$ одну решётку? Найдите
матрицу перехода и её определитель. Как изменятся $\det \Lat$ и $\mathcal H$, если все векторы
базиса умножить на $c > 0$?}

\section{Задачи SVP и CVP}

\begin{definition}
Пусть $\Lat \subset \RR^n$ "--- решётка.
\begin{enumerate}
  \item \emph{Задача о кратчайшем векторе} (SVP, shortest vector problem): найти ненулевой
    $v \in \Lat$ минимальной длины. Его длина обозначается $\lambda_1(\Lat)$.
  \item \emph{Задача о ближайшем векторе} (CVP, closest vector problem): по вектору
    $w \in \RR^n$ (не обязательно из $\Lat$) найти $v \in \Lat$, минимизирующий $\norm{w - v}$.
  \item \emph{Приближённые варианты} с множителем $\psi(n) \ge 1$: apprSVP "--- найти ненулевой
    $v \in \Lat$ с $\norm{v} \le \psi(n) \lambda_1(\Lat)$; apprCVP "--- найти $v \in \Lat$ с
    $\norm{w - v} \le \psi(n) \min_{u \in \Lat}\norm{w - u}$.
\end{enumerate}
\end{definition}

Кратчайший вектор существует (решётка дискретна), но не единствен: вместе с $v$ кратчайшим
является и $-v$, а бывают и другие (в $\ZZ^2$ их четыре). Решения задач даны в терминах
решётки, а на вход алгоритму подаётся \emph{базис}; трудность задачи сильно зависит от того,
насколько хорош этот базис.

\paragraph{Сложность} (без доказательства, см.~\cite[гл.~3, 4]{micciancio-goldwasser}).
\begin{itemize}
  \item CVP NP-трудна (П.~ван Эмде Боас, 1981).
  \item SVP NP-трудна относительно вероятностных сведений (М.~Айтаи, 1998).
  \item Приближённые варианты остаются трудными при постоянных множителях $\psi$, тогда как
    при $\psi(n) = 2^{n/2}$ они решаются за полиномиальное время алгоритмом LLL (лекция~12).
    Криптосистемы на решётках используют множители, полиномиальные по $n$, "--- для них
    неизвестны ни NP-трудность, ни эффективные (в том числе квантовые) алгоритмы.
\end{itemize}

На практике CVP обычно не проще SVP: известные алгоритмы решают обе задачи за время,
экспоненциальное по $n$, и при $n$ порядка нескольких сотен они становятся неприменимы.

\subsection{Теоремы Эрмита и Минковского}

Сколь коротким может быть кратчайший вектор решётки с данным определителем?

\begin{theorem}[теорема Минковского о выпуклом теле]\label{l11:minkowski}
Пусть $\Lat \subset \RR^n$ "--- решётка полного ранга, $S \subset \RR^n$ "--- ограниченное
измеримое выпуклое множество, симметричное относительно нуля ($S = -S$). Если
$\operatorname{vol}(S) > 2^n \det \Lat$, то $S$ содержит ненулевой вектор решётки. Если $S$ к тому
же замкнуто, достаточно условия $\operatorname{vol}(S) \ge 2^n \det \Lat$.
\end{theorem}

\begin{proof}
Пусть $\mathcal F$ "--- фундаментальная область некоторого базиса, и $S' = \frac12 S$, так что
$\operatorname{vol}(S') = 2^{-n}\operatorname{vol}(S) > \det \Lat = \operatorname{vol}(\mathcal F)$.
По утверждению~\ref{l11:tiling} множество $S'$ разбивается на непересекающиеся куски
$S'_v = S' \cap (\mathcal F + v)$, $v \in \Lat$ (их конечное число, так как $S'$ ограничено), и
сдвинутые куски $S'_v - v$ лежат в $\mathcal F$. Если бы они попарно не пересекались, то
\[
  \operatorname{vol}(S') = \sum_v \operatorname{vol}(S'_v) = \sum_v \operatorname{vol}(S'_v - v)
  \le \operatorname{vol}(\mathcal F),
\]
противоречие. Значит, $x_1 - v_1 = x_2 - v_2$ для некоторых $x_1, x_2 \in S'$ и $v_1 \ne v_2$.
Тогда $0 \ne v_1 - v_2 = x_1 - x_2 = \frac12(2x_1) + \frac12(-2x_2)$, где $2x_1 \in S$ и
$-2x_2 \in S$ (симметрия), и по выпуклости $v_1 - v_2 \in S$.

Для замкнутого $S$ применим доказанное к $(1 + \varepsilon) S$, $\varepsilon = 1/k$: каждое из
этих множеств содержит ненулевую точку решётки, все они лежат в ограниченном множестве $2S$,
где точек решётки конечное число. Значит, одна и та же точка $v \ne 0$ лежит в
$(1 + 1/k)S$ при бесконечно многих $k$, а так как $S$ замкнуто, $v \in S$.
\end{proof}

\begin{theorem}[Эрмит]\label{l11:hermite}
В каждой решётке $\Lat \subset \RR^n$ полного ранга есть ненулевой вектор $v$ с
\[
  \norm{v} \le \sqrt{n}\,(\det \Lat)^{1/n}.
\]
\end{theorem}

\begin{proof}
Пусть $a = (\det \Lat)^{1/n}$, $S = [-a, a]^n$ "--- замкнутый симметричный выпуклый куб,
$\operatorname{vol}(S) = (2a)^n = 2^n \det \Lat$. По теореме~\ref{l11:minkowski} в $S$ есть ненулевой
$v \in \Lat$, и $\norm{v} \le \sqrt{n a^2} = \sqrt n\, a$.
\end{proof}

\begin{remark}
Наименьшая константа $\gamma_n$ с $\lambda_1(\Lat)^2 \le \gamma_n (\det \Lat)^{2/n}$ для всех
решёток размерности $n$ называется \emph{константой Эрмита}. Теорема~\ref{l11:hermite} даёт
$\gamma_n \le n$; применяя теорему Минковского к шару, получают
$\gamma_n \le \frac{4}{\pi}\Gamma(1 + \frac n2)^{2/n} \approx \frac{2n}{\pi e}$. Точные значения
известны лишь для $n \le 8$ и $n = 24$; например, $\gamma_2 = 2/\sqrt3$
(следствие~\ref{l11:gamma2} ниже).
\end{remark}

\paragraph{Гауссова эвристика.} Для «случайной» решётки число её точек в шаре $B_R$ радиуса $R$
приблизительно равно $\operatorname{vol}(B_R)/\det \Lat$. Приравнивая это к $1$, получаем
ожидаемую длину кратчайшего вектора
\[
  \sigma(\Lat) = \left(\frac{\Gamma(1 + n/2)\det \Lat}{\pi^{n/2}}\right)^{1/n} \approx
  \sqrt{\frac{n}{2\pi e}}\,(\det \Lat)^{1/n} \qquad (n \to \infty).
\]
Это не теорема, а эвристика, хорошо подтверждаемая экспериментами для «типичных» решёток.
Если в решётке есть вектор, \emph{намного} более короткий, чем $\sigma(\Lat)$, он, как
правило, единственный такой (с точностью до знака) и его легче найти алгоритмами редукции
"--- на этом основаны многие атаки. В примере~\ref{l11:congr-example} $n = 2$,
$\det \Lat = q$ и $\sigma(\Lat) = \sqrt{q/\pi} \approx 564$; кратчайший вектор "--- ключ
$(325, 534)$ длины $\approx 625$ (оценка Эрмита: $\sqrt{2q} \approx 1414$), тогда как векторы
исходного базиса $(1, h), (0, q)$ имеют длину порядка $q = 10^6$.

\checkyourself{Почему теорема Минковского неверна без условия симметрии $S = -S$? Без
выпуклости? Приведите примеры на плоскости для решётки $\ZZ^2$.}

\section{Алгоритм Гаусса}\label{l11:sec-gauss}

В размерности $2$ задача SVP решается быстро алгоритмом, восходящим к Лагранжу и Гауссу.
Идея "--- как в алгоритме Евклида: вычитаем из длинного вектора подходящее кратное короткого,
пока это уменьшает длину.

\begin{algorithm}[Гаусс, редукция базиса размерности $2$]\label{l11:alg-gauss}
\begin{algorithmic}[1]
\Require базис $v_1, v_2$ решётки $\Lat \subset \RR^2$.
\Ensure базис $v_1, v_2$ того же $\Lat$, где $v_1$ "--- кратчайший вектор $\Lat$.
\Repeat
  \If{$\norm{v_2} < \norm{v_1}$} \State поменять $v_1 \leftrightarrow v_2$ \EndIf
  \State $m \gets \round{v_1 \cdot v_2 / \norm{v_1}^2}$ \Comment{ближайшее целое}
  \State $v_2 \gets v_2 - m v_1$
\Until{$m = 0$}
\State \Return $(v_1, v_2)$
\end{algorithmic}
\end{algorithm}

Вектор $v_2 - \mu v_1$ при $\mu = v_1 \cdot v_2/\norm{v_1}^2$ "--- проекция $v_2$ на прямую,
ортогональную $v_1$ (это $v_2^*$), т.\,е. кратчайший из векторов $v_2 - t v_1$, $t \in \RR$;
оставаясь в решётке, берём ближайшее целое $t = m$. Смена базиса на каждом шаге задаётся
матрицей $\begin{psmallmatrix}1 & 0\\ -m & 1\end{psmallmatrix}$ или
$\begin{psmallmatrix}0 & 1\\ 1 & 0\end{psmallmatrix}$ из $\mathrm{GL}_2(\ZZ)$, поэтому решётка не
меняется (теорема~\ref{l11:basis-change}).

\begin{theorem}\label{l11:gauss}
Алгоритм~\ref{l11:alg-gauss} останавливается. Для результата $(v_1, v_2)$ выполнено
$\norm{v_1} \le \norm{v_2}$, $\abs{v_1 \cdot v_2} \le \frac12 \norm{v_1}^2$, и
\begin{enumerate}
  \item $v_1$ "--- кратчайший ненулевой вектор $\Lat$: $\norm{v_1} = \lambda_1(\Lat)$;
  \item $v_2$ "--- кратчайший среди векторов $\Lat$, не кратных $v_1$.
\end{enumerate}
\end{theorem}

\begin{proof}
\emph{Остановка.} Пусть в некоторой итерации (после возможной перестановки)
$\norm{v_1} \le \norm{v_2}$, $\mu = v_1 \cdot v_2 / \norm{v_1}^2$ и $m = \round{\mu} \ne 0$.
Новый вектор $v_2' = v_2 - m v_1$ имеет $v_1 \cdot v_2' / \norm{v_1}^2 = \mu - m \in
[-\frac12, \frac12)$. Если в следующей итерации перестановки нет, то $m = \round{\mu - m} = 0$
и алгоритм останавливается. Если же перестановка есть, то $\norm{v_2'} < \norm{v_1}$, т.\,е.
первый вектор базиса \emph{строго} укорачивается. Итак, до остановки $\norm{v_1}$ строго
убывает от итерации к итерации, а все $v_1$ "--- векторы решётки длины не больше исходных;
таких векторов конечное число (решётка дискретна), поэтому алгоритм останавливается.

\emph{Свойства результата.} При остановке перестановки нет, т.\,е. $\norm{v_1} \le \norm{v_2}$, и
$\round{\mu} = 0$, т.\,е. $\abs{\mu} \le \frac12$, что равносильно $\abs{v_1 \cdot v_2} \le
\frac12\norm{v_1}^2$. Пусть $v = a_1 v_1 + a_2 v_2$, $(a_1, a_2) \in \ZZ^2 \setminus \{0\}$. Тогда
\[
  \norm{v}^2 = a_1^2 \norm{v_1}^2 + 2a_1a_2\, v_1 \cdot v_2 + a_2^2 \norm{v_2}^2
  \ge a_1^2 \norm{v_1}^2 - \abs{a_1a_2}\norm{v_1}^2 + a_2^2 \norm{v_2}^2.
\]
(1) Так как $\norm{v_2} \ge \norm{v_1}$, правая часть не меньше
$(a_1^2 - \abs{a_1a_2} + a_2^2)\norm{v_1}^2 \ge \norm{v_1}^2$: для целых $a_1, a_2$, не равных
одновременно нулю, $a_1^2 - \abs{a_1a_2} + a_2^2 \ge 1$ (это целое число, и
$a_1^2 - \abs{a_1a_2} + a_2^2 = (\abs{a_1} - \frac12\abs{a_2})^2 + \frac34 a_2^2 > 0$).

(2) Пусть $a_2 \ne 0$. Если $\abs{a_1} \ge \abs{a_2}$, то $a_1^2 - \abs{a_1a_2} \ge 0$ и правая
часть не меньше $a_2^2 \norm{v_2}^2 \ge \norm{v_2}^2$. Если $\abs{a_1} < \abs{a_2}$, то
$a_1^2 - \abs{a_1a_2} < 0$, и, заменяя $\norm{v_1}^2$ на большее $\norm{v_2}^2$, получаем
$\norm{v}^2 \ge (a_1^2 - \abs{a_1a_2} + a_2^2)\norm{v_2}^2 \ge \norm{v_2}^2$.
\end{proof}

\begin{corollary}\label{l11:gamma2}
Для любой решётки $\Lat \subset \RR^2$: $\lambda_1(\Lat)^2 \le \frac{2}{\sqrt3}\det \Lat$, т.\,е.
$\gamma_2 \le 2/\sqrt3$ (равенство "--- для «шестиугольной» решётки $\Lat((1, 0), (\frac12,
\frac{\sqrt3}{2}))$).
\end{corollary}

\begin{proof}
Для результата алгоритма Гаусса
$\det \Lat = \norm{v_1}\norm{v_2^*}$ (теорема~\ref{l11:hadamard}), где
$\norm{v_2^*}^2 = \norm{v_2}^2 - \mu^2 \norm{v_1}^2 \ge \norm{v_1}^2 (1 - \frac14)$. Значит,
$(\det \Lat)^2 \ge \frac34 \norm{v_1}^4 = \frac34 \lambda_1^4$.
\end{proof}

\paragraph{Сложность.} Каждая итерация "--- $\bigO(1)$ операций над числами длины входа.
Можно показать, что число итераций есть $\bigO(\log \max \norm{v_i})$ (как у алгоритма Евклида,
частным случаем которого он является: для $v_1 = (a, 0)$, $v_2 = (b, 0)$ в «решётке» на прямой
алгоритм Гаусса превращается в алгоритм Евклида с округлением до ближайшего)
\cite[разд.~17.1]{galbraith}.

\begin{example}\label{l11:gauss-hand}
$v_1 = (16, 26)$, $v_2 = (20, 34)$, $\det \Lat = 16 \cdot 34 - 26 \cdot 20 = 24$,
$\mathcal H \approx 0{,}141$.
\[
\begin{array}{c|r|r|r|r|r}
 & v_1 & v_2 & \norm{v_1}^2 & v_1 \cdot v_2 & m \\ \hline
1 & (16, 26) & (20, 34) & 932 & 1204 & 1 \\
2^{\leftrightarrow} & (4, 8) & (16, 26) & 80 & 272 & 3 \\
3^{\leftrightarrow} & (4, 2) & (4, 8) & 20 & 32 & 2 \\
4 & (4, 2) & (-4, 4) & 20 & -8 & 0
\end{array}
\]
(${}^{\leftrightarrow}$ "--- была перестановка). Итог: $\lambda_1 = \norm{(4, 2)} = \sqrt{20}$,
второй вектор $(-4, 4)$; новый базис почти ортогонален: $\sqrt{20}\cdot\sqrt{32} \approx 25{,}3$
против $\det \Lat = 24$, $\mathcal H \approx 0{,}974$.
\end{example}

\taskexample{lattice-gauss}

\paragraph{Атака на конгруэнтную криптосистему.} Применим алгоритм Гаусса к базису
$(1, h), (0, q)$ примера~\ref{l11:congr-example}, $h = 529\,234$, $q = 1\,000\,003$:
\[
\begin{array}{c|r|r|r}
 & v_1 & v_2 & m \\ \hline
1 & (1, 529\,234) & (0, 1\,000\,003) & 2 \\
2^{\leftrightarrow} & (-2, -58\,465) & (1, 529\,234) & -9 \\
3^{\leftrightarrow} & (-17, 3\,049) & (-2, -58\,465) & -19 \\
4^{\leftrightarrow} & (-325, -534) & (-17, 3\,049) & -4 \\
5 & (-325, -534) & (-1\,317, 913) & 0
\end{array}
\]
Кратчайший вектор $\pm(325, 534) = (f, g)$ "--- закрытый ключ Алисы. Отметим, что Еве не
обязательно найти именно $(f, g)$: годится любой короткий вектор $(f', g') \in \Lat$, с которым
расшифрование корректно (доказательство утверждения~\ref{l11:congr-correct} использует лишь
$f'h \equiv g'$ и малость $f', g'$).

\taskexample{lattice-congruential}

\checkyourself{Что вернёт алгоритм Гаусса, если исходный базис уже удовлетворяет условиям
$\norm{v_1} \le \norm{v_2}$, $\abs{v_1 \cdot v_2} \le \frac12\norm{v_1}^2$? Покажите, что для
такого базиса угол между $v_1$ и $v_2$ лежит в $[60^\circ, 120^\circ]$.}

\section{Алгоритм Бабаи}

Пусть дан базис $v_1, \dots, v_n$ решётки и вектор $w \in \RR^n$. Если базис ортогонален,
CVP решается легко: запишем $w = t_1 v_1 + \cdots + t_n v_n$ и округлим коэффициенты.
Алгоритм Бабаи делает то же для произвольного базиса \cite{babai}.

\begin{algorithm}[Бабаи, округление]\label{l11:alg-babai}
\begin{algorithmic}[1]
\Require базис $B = (v_1, \dots, v_n)$ решётки $\Lat \subset \RR^n$, вектор $w \in \RR^n$.
\Ensure вектор $v \in \Lat$, близкий к $w$.
\State решить систему $w = t_1 v_1 + \cdots + t_n v_n$, т.\,е. $t = wB^{-1}$
\State $a_i \gets \round{t_i}$, $i = 1, \dots, n$
\State \Return $v = a_1 v_1 + \cdots + a_n v_n$
\end{algorithmic}
\end{algorithm}

Сложность "--- $\bigO(n^3)$ арифметических операций (решение линейной системы). В размерности
$2$ коэффициенты удобно находить по правилу Крамера:
$t_1 = \det\begin{psmallmatrix} w \\ v_2\end{psmallmatrix} / \det B$,
$t_2 = \det\begin{psmallmatrix} v_1 \\ w\end{psmallmatrix} / \det B$.

\begin{proposition}\label{l11:babai-bound}
Пусть $v$ "--- результат алгоритма~\ref{l11:alg-babai}. Тогда
\begin{enumerate}
  \item $\norm{w - v} \le \frac12 (\norm{v_1} + \cdots + \norm{v_n})$;
  \item если базис ортогонален, то $v$ "--- ближайший к $w$ вектор решётки.
\end{enumerate}
\end{proposition}

\begin{proof}
(1) $w - v = \sum (t_i - a_i) v_i$, где $\abs{t_i - a_i} \le \frac12$; применяем неравенство
треугольника.

(2) Для любого $u = \sum b_i v_i \in \Lat$ по теореме Пифагора
$\norm{w - u}^2 = \sum (t_i - b_i)^2\norm{v_i}^2$. Каждое слагаемое минимально при $b_i =
\round{t_i} = a_i$, поэтому минимум достигается при $u = v$.
\end{proof}

Для почти ортогонального базиса (с $\mathcal H(B)$, близким к $1$) алгоритм Бабаи, как правило,
тоже находит ближайший вектор, а для плохого базиса оценка (1) велика и результат может быть
далёк от ответа.

\begin{example}\label{l11:babai-hand}
Решётка рис.~\ref{l11:fig-bases}, $w = (-\frac{11}{4}, \frac14)$.

\emph{Хороший базис} $v_1 = (2, 1)$, $v_2 = (-1, 1)$: $2t_1 - t_2 = -\frac{11}4$,
$t_1 + t_2 = \frac14$, откуда $t_1 = -\frac56$, $t_2 = \frac{13}{12}$. Округляем:
$v = -v_1 + v_2 = (-3, 0)$, $\norm{w - v} = \sqrt{\frac1{16} + \frac1{16}} \approx 0{,}354$; это
ближайший вектор решётки.

\emph{Плохой базис} $u_1 = (-4, 1) = -v_1 + 2v_2$, $u_2 = (-5, 2) = -v_1 + 3v_2$:
$w = s_1 u_1 + s_2 u_2$ при $s_1 = \frac{17}{12} \approx 1{,}42$, $s_2 = -\frac{7}{12} \approx -0{,}58$;
округляем: $u_1 - u_2 = (1, -1)$, $\norm{w - (1, -1)} = \sqrt{\frac{225}{16} + \frac{25}{16}} \approx
3{,}953$ "--- в одиннадцать раз дальше. Оценки (1): $\frac12(\sqrt5 + \sqrt2) \approx 1{,}83$ и
$\frac12(\sqrt{17} + \sqrt{29}) \approx 4{,}75$.
\end{example}

\taskexample{lattice-babai-good}

Если базис плохой, его сначала улучшают (в размерности $2$ "--- алгоритмом Гаусса, в общем
случае "--- алгоритмом LLL, лекция~12), а затем применяют алгоритм Бабаи:

\taskexample{lattice-babai-bad}

\taskexample{lattice-babai-compare}

\subsection{Алгоритм ближайшей плоскости}

Второй алгоритм Бабаи использует ортогонализацию Грама"--~Шмидта и даёт гарантию, зависящую
только от длин $\norm{v_i^*}$. Решётка $\Lat$ лежит в объединении параллельных гиперплоскостей
$c\,v_n + \Span(v_1, \dots, v_{n-1})$, $c \in \ZZ$; выбираем ближайшую к $w$ из них, проецируем
задачу на неё и повторяем в размерности $n - 1$.

\begin{algorithm}[Бабаи, ближайшая плоскость]\label{l11:alg-nearest-plane}
\begin{algorithmic}[1]
\Require базис $v_1, \dots, v_n$ решётки $\Lat$, вектор $w$.
\Ensure вектор $v \in \Lat$, близкий к $w$.
\State вычислить $v_1^*, \dots, v_n^*$ (Грам"--~Шмидт); $w' \gets w$
\For{$i = n, n-1, \dots, 1$}
  \State $c_i \gets \round{w' \cdot v_i^* / \norm{v_i^*}^2}$
  \State $w' \gets w' - c_i v_i$
\EndFor
\State \Return $v = w - w' = c_1 v_1 + \cdots + c_n v_n$
\end{algorithmic}
\end{algorithm}

\begin{proposition}\label{l11:nearest-plane}
Для результата алгоритма~\ref{l11:alg-nearest-plane}
\begin{enumerate}
  \item $w - v = \sum_{i=1}^n \varepsilon_i v_i^*$ с $\abs{\varepsilon_i} \le \frac12$, поэтому
    $\norm{w - v}^2 \le \frac14 \sum_{i=1}^n \norm{v_i^*}^2$;
  \item если $\norm{w - u} < \frac12 \min_i \norm{v_i^*}$ для некоторого $u \in \Lat$, то $v = u$
    (алгоритм находит ближайший вектор).
\end{enumerate}
\end{proposition}

\begin{proof}
Запишем $w' = \sum_j \kappa_j v_j^*$. Вычитание $c_i v_i$ не меняет коэффициенты $\kappa_j$ при
$j > i$ (так как $v_i \in \Span(v_1^*, \dots, v_i^*) \perp v_j^*$) и уменьшает $\kappa_i$ ровно на
$c_i$ (так как $v_i \cdot v_i^* = \norm{v_i^*}^2$). На шаге $i$ имеем $\kappa_i = w' \cdot v_i^* /
\norm{v_i^*}^2$, и после вычитания $\kappa_i - \round{\kappa_i} \in [-\frac12, \frac12]$; более
поздние шаги ($j < i$) его не меняют. Отсюда (1) по теореме Пифагора.

(2) Пусть $u = \sum a_i v_i$, $w - u = \sum \delta_i v_i^*$, где
$\abs{\delta_i} \le \norm{w - u}/\norm{v_i^*} < \frac12$. Индукцией по $i = n, \dots, 1$ покажем
$c_i = a_i$. Пусть $c_j = a_j$ при $j > i$; тогда перед шагом $i$
$w' = w - \sum_{j > i} a_j v_j = (w - u) + \sum_{j \le i} a_j v_j$, и коэффициент при $v_i^*$ равен
$\delta_i + a_i$ (векторы $v_j$, $j < i$, не дают вклада, у $v_i$ он равен $1$). Так как
$\abs{\delta_i} < \frac12$, $c_i = \round{a_i + \delta_i} = a_i$.
\end{proof}

Для базиса, редуцированного алгоритмом LLL, длины $\norm{v_i^*}$ убывают не слишком быстро,
и Бабаи доказал, что тогда $\norm{w - v} \le 2^{n/2} \min_{u \in \Lat}\norm{w - u}$, т.\,е.
apprCVP решается с множителем $2^{n/2}$ (без доказательства; лекция~12, \cite{babai}).

\checkyourself{Покажите, что в размерности $2$ для ортогонального базиса алгоритмы округления и
ближайшей плоскости совпадают. Приведите пример, когда они дают разные ответы.}

\section{Криптосистемы на решётках: обзор}

\begin{remark}[GGH]
В криптосистеме Гольдрайха"--~Гольдвассер"--~Халеви (1997) закрытый ключ "--- хороший базис
$V$ решётки, открытый "--- плохой базис $W = UV$ той же решётки. Сообщение "--- целочисленный
вектор $m$; шифртекст $e = mW + r$, где $r$ "--- малый случайный вектор. Алиса применяет
алгоритм Бабаи с хорошим базисом $V$, находит ближайший к $e$ вектор $mW$ и $m = (mW)W^{-1}$.
Ева с плохим базисом $W$ решает трудную задачу CVP. Конкретные параметры GGH были взломаны
П.~Нгуеном (1999), но сама идея «хороший базис "--- секрет, плохой "--- открыт» лежит в основе
многих схем, в том числе подписей \cite[разд.~7.8]{hps}.
\end{remark}

\begin{remark}[рюкзак и NTRU]
Рюкзачная криптосистема Меркла"--~Хеллмана (1978) основана на задаче о сумме подмножества:
по $M_1, \dots, M_n$ и $S$ найти $x_i \in \{0, 1\}$ с $\sum x_i M_i = S$. Решение даёт
короткий вектор $(2x_1 - 1, \dots, 2x_n - 1, 0)$ решётки со строками $(2e_i, M_i)$,
$i = 1, \dots, n$, и $(1, \dots, 1, S)$; при больших $n$ его находит алгоритм LLL
(Лагариас"--~Одлыжко, 1985), и рюкзачные системы взломаны \cite[разд.~7.2]{hps}.
Криптосистема NTRU (Хоффштейн"--~Пайфер"--~Сильверман, 1996) обобщает конгруэнтную
криптосистему: числа $f, g, h$ заменяются многочленами степени меньше $N$, а решётка
$\Lat((1, h), (0, q))$ "--- решёткой размерности $2N$ с базисом
$\begin{psmallmatrix} I & H \\ 0 & qI \end{psmallmatrix}$, где $H$ "--- циркулянтная матрица
многочлена $h$ \cite[разд.~7.10--7.11]{hps}. При $N$ в несколько сотен кратчайший вектор
такой решётки найти не удаётся.
\end{remark}

\checkyourself{Почему в GGH нельзя публиковать хороший базис? Почему недостаточно держать
в секрете лишь матрицу $U$, если $V$ общеизвестен?}

\section{Что дальше}

\begin{itemize}
  \item В размерности $n > 2$ точного аналога алгоритма Гаусса с полиномиальной сложностью нет.
    Алгоритм LLL (Ленстра"--~Ленстра"--~Ловас, 1982, \cite{lll}) за полиномиальное время находит
    базис с $\norm{v_1} \le 2^{(n-1)/2}\lambda_1(\Lat)$ "--- лекция~12.
  \item Комбинация LLL и алгоритма ближайшей плоскости решает apprCVP; на практике
    используются более сильные редукции (BKZ) и перебор.
  \item Современные постквантовые схемы (Kyber, Dilithium, Falcon) основаны на задачах LWE и
    NTRU, сводящихся к apprSVP с полиномиальным множителем.
\end{itemize}

\checkyourselfsection

\begin{furtherreading}
  \readitem[разд.~7.1]{hps}{конгруэнтная криптосистема и атака на неё}
  \readitem[разд.~7.3--7.6]{hps}{решётки, фундаментальная область, неравенство Адамара,
    SVP и CVP, теоремы Эрмита и Минковского, гауссова эвристика, алгоритм Бабаи}
  \readitem[разд.~7.8, 7.10--7.11, 7.13.1]{hps}{GGH, NTRU; алгоритм Гаусса}
  \readitem[гл.~16--18]{galbraith}{решётки, редукция в размерности $2$ (разд.~17.1),
    алгоритмы Бабаи для CVP}
  \readitem[гл.~1--4]{micciancio-goldwasser}{теория решёток, сложность SVP и CVP}
  \readitem[гл.~2]{nguyen-vallee}{константа Эрмита и алгоритм Лагранжа"--~Гаусса}
  \readitem{babai}{оригинальная статья: алгоритмы округления и ближайшей плоскости}
  \readitem{regev-lattices}{конспект курса О.~Регева (лекция~1: решётки, определитель,
    теорема Минковского)}
\end{furtherreading}
\end{lecture}
